
% -----------------------------------------------------------
\section{Search Becomes Abundant}

Something genuinely unprecedented has happened to the economics of exploration.

For centuries, the cost of examining an idea was proportional to the cost of
developing it. A scientist wishing to compare three competing hypotheses had
to run three experiments. An architect wishing to evaluate alternative structural
systems had to produce three sets of drawings. A programmer wishing to explore
three approaches to a design problem had to implement three programs. Exploration
was bounded by implementation capacity, and implementation capacity was scarce.
The practical consequence was that inquiry was forced to narrow before it widened.
Exploration had to be preceded by sufficient prior reasoning to justify the
expenditure of implementation effort. The thinker compressed the search space
mentally before exploring it empirically.

Recent systems capable of generating candidate solutions at high speed and low
marginal cost have substantially altered this relationship. The collapse of
generation costs means that a single individual can now explore regions of
possibility space that would previously have required teams of specialists or
months of concentrated labor. Three candidate architectures can be sketched in
hours rather than weeks. Ten alternative formulations of a problem can be
examined before committing to any one of them. The cost that previously
compelled premature convergence has been dramatically reduced.

This is a genuine enlargement of human cognitive capacity. It is worth dwelling
on the magnitude of the change before examining its consequences. The ability
to defer commitment, to explore widely before narrowing, to hold multiple
candidate structures in play simultaneously --- these are not trivial affordances.
Scientific method in practice has always been constrained by the cost of
generating and testing hypotheses. Design practice has always been constrained
by the cost of prototyping. The relaxation of those constraints, even partially,
represents a qualitative shift in how inquiry can proceed.

The question is what this shift implies for the overall structure of productive work.

A common answer is that increased generation capacity translates directly into
increased productive output. If generation is faster, more can be produced in
a given time. This answer is not wrong, but it is incomplete in a way that
matters. It focuses on what the expansion of search capacity makes easier while
neglecting what it leaves unchanged --- and what it makes harder.

% -----------------------------------------------------------
\section{The Asymmetry Between Search and Orientation}

To identify what the expansion of search leaves unchanged, it is useful to
distinguish two components of any goal-directed inquiry.

The first component is \emph{search}: the process of generating candidate states,
solutions, hypotheses, or structures. Search explores a possibility space. It
asks, in various ways, what could exist here. The outputs of search are
candidates --- proposals that may or may not satisfy the objectives of the inquiry.

The second component is \emph{orientation}: the process of maintaining a coherent
objective through time, evaluating candidates against that objective, determining
which regions of possibility space are worth exploring, and recognizing when
the inquiry has made genuine progress as opposed to local movement. Orientation
asks, in various ways, what is worth reaching and whether one is still moving
toward it.

These two components are not symmetric in how they respond to changes in
generation capacity.

Search scales readily. When generation is cheap, more candidates can be produced.
The range of positions, architectures, solutions, and alternatives available for
examination expands. The breadth of exploration increases. More of possibility
space becomes accessible per unit of time.

Orientation does not scale through the same mechanisms as search. It is worth
being precise about why, because orientation is not a single capacity. It
comprises at least four distinguishable functions. The first is \emph{evaluation}:
determining whether a candidate satisfies some criterion. The second is
\emph{goal maintenance}: keeping the original objective in view across an
extended sequence of local decisions. The third is \emph{trajectory assessment}:
judging whether the current path, not merely the current state, is moving
toward the objective. The fourth is \emph{goal revision}: recognizing when
the original objective was itself mistaken or incomplete, and adjusting it
without losing continuity with what was genuinely intended.

These four functions scale differently. Evaluation, when the criterion is
well-specified and computable, can often be automated or parallelized. It is
the component of orientation most amenable to computational assistance. Goal
maintenance is harder to delegate, because it requires a representation of
the objective that is not itself a product of the local context. Trajectory
assessment is harder still, because it requires comparing sequences rather
than states. Goal revision is the least tractable of all, because it requires
the kind of value-laden judgment that cannot be reduced to optimization against
a fixed criterion --- it is, in a sense, the meta-orientation that governs
all the others.

The characteristic error of an age of abundant search is to treat evaluation
as though it were the whole of orientation, and to conclude that because
evaluation can scale, orientation can scale. The higher levels of orientation
--- goal maintenance, trajectory assessment, and especially goal revision ---
remain substantially tied to the cognitive and institutional resources of the
inquiring agent. They require continuity of judgment that cannot be easily
parallelized or offloaded.

\begin{principle}
Search and orientation scale at different rates. Search scales computationally,
through parallelism and faster generation. Orientation scales institutionally
and cognitively, through review, shared constraints, memory systems, and the
continuity of judgment. Confusing these scaling regimes is the characteristic
error of an age of abundant search.
\end{principle}

\begin{principle}
As search becomes abundant, orientation becomes scarce.
\end{principle}

\subsection*{A Concrete Example}

The distinction between search and orientation is not unique to automated
generation systems. Similar transitions appear whenever technological advances
expand the set of reachable states faster than the capacity to evaluate them.

A recent example comes from dense biomedical image registration. Jena,
Chaudhari, and Gee describe a registration framework that dramatically reduces
the computational cost of aligning large volumetric datasets across modalities
and species~\cite{fireants2026}. Hyperparameter studies that would previously
have required years of computation become feasible in hours, enabling researchers
to explore hundreds of candidate configurations before committing to any single
registration strategy.

The significance of this result is not merely speed. The enlarged computational
budget expands the region of parameter space that can be examined before
commitment. More candidate solutions become reachable. Yet the scientific
problem remains unchanged. Researchers must still determine which metrics
matter, which regions of parameter space correspond to biologically meaningful
solutions, and when exploration has been sufficient to justify formalization.
The search space expands. The requirement for orientation does not. Indeed,
the need for orientation becomes more acute precisely because more possibilities
can now be examined without the natural constraint that expense formerly imposed.

The practical implication of this asymmetry follows directly. When both search
and orientation were expensive, the limiting resource was search. Productive
work was constrained primarily by the ability to generate candidates. But as
search becomes abundant, the asymmetry becomes visible. Orientation, which was
never the bottleneck when search was expensive, becomes the bottleneck when
search becomes cheap. The limiting factor in productive inquiry shifts from
the capacity to explore to the capacity to remain oriented within an increasingly
large explored space.

This shift has not been widely recognized, for a simple reason: the expansion
of search capacity is highly visible, while the constraint imposed by orientation
is largely invisible. One can count candidates generated. One cannot easily count
orientations maintained. The measurable thing improves dramatically; the
unmeasured thing remains constant; and the temptation is to interpret the
improvement in the measurable thing as an overall increase in productive capacity.
But what has increased is search. What has not increased is the capacity to
make use of search while remaining coherently aimed.

% -----------------------------------------------------------
\section{Exploration Before Formalization}

The asymmetry between search and orientation has direct implications for how
productive inquiry should be structured.

A common model of goal-directed work runs as follows. The agent specifies the
objective in advance. The specification is then handed to a process --- human
or automated --- that implements it. The output is evaluated and refined.
This model has the virtue of clarity. The objective is explicit. The
implementation follows from the specification.

The difficulty with this model is that it assumes the agent can produce a
complete and accurate specification before exploration has occurred. In practice,
specifications are hypotheses about what is wanted. They may be partially correct.
They will almost certainly be incomplete. They will contain assumptions that
appear reasonable before exploration but reveal themselves as mistaken once
alternatives are examined. The specification encodes the agent's prior beliefs
about the solution space. Prior beliefs about solution spaces are often wrong
in ways that can only be discovered by exploring the space.

Productive inquiry, whether in science, design, engineering, or any other domain
requiring the navigation of a large possibility space, characteristically follows
a different structure. Exploration precedes formalization. The agent examines
the territory before committing to an architecture. Prototypes are built before
interfaces are specified. Experiments are run before theories are committed to
paper. The laboratory prototype is not a defective version of the final instrument.
It is a device for exposing structure that the inquirer does not yet have access
to. Once the structure has been exposed, formalization can proceed from genuine
understanding rather than from prior belief.

\begin{observation}
Formalization is an irreversible operation. Every abstraction eliminates
possibilities. Every interface commitment closes off alternative structures.
Every architectural decision forecloses some future modifications. The cost
of premature formalization is therefore not merely that the resulting structure
is suboptimal, but that it is difficult to revise. Premature commitment is
not an implementation error but an orientation error: it eliminates regions
of possibility space before their value has been established.
\end{observation}

The expansion of search capacity makes the separation between exploration and
formalization more important, not less. When generation is cheap, the quantity
of candidates that can be produced before commitment is dramatically larger.
The territory that can be explored before formalization is much wider. But
exploiting this expanded exploration capacity requires maintaining the distinction.
If cheap generation is used to implement specifications more quickly rather
than to explore more widely before committing to specifications, the new
capacity is being used to do the old thing faster rather than to do something
qualitatively different. The result is more rapid convergence on a solution
that may be no better chosen than before --- merely more quickly arrived at.

The inquiry that benefits most from abundant search is inquiry that defers
formalization until exploration has been sufficient to justify commitment.
The exploration phase exposes structure. The formalization phase stabilizes it.
Collapsing them eliminates the epistemic benefit of cheap generation.

The biomedical registration case is instructive here as well. The authors
emphasize that optimal registration parameters are not universal: they differ
across modalities, species, and experimental objectives~\cite{fireants2026}.
Faster computation does not eliminate the need to determine which parameters
are appropriate for which contexts. It makes the exploration of that question
tractable for the first time. The system does not assume the correct parameters
in advance and implement them efficiently. It expands the space that can be
explored before any commitment is made. That is precisely the structure this
section is describing: abundant search enabling genuine exploration, rather
than more rapid convergence on a prior assumption.

% -----------------------------------------------------------
\section{Goal Continuity}

A further consequence of the asymmetry between search and orientation concerns
what might be called goal continuity: the preservation of the agent's original
objective through an extended sequence of local decisions.

In any iterative process, each step produces an intermediate state. The agent
evaluates the intermediate state, makes a decision about the next step, and
proceeds. In a short process, with few steps, the original objective remains
clearly in view throughout. But in a long iterative process, the relationship
between each local step and the original objective can become increasingly
indirect. The agent is evaluating the intermediate state, not the objective
directly. The decision about the next step is made relative to the intermediate
state. If the intermediate state is an imperfect proxy for the objective ---
if it captures some aspects of what was originally wanted while losing others
--- then each successive step may optimize the proxy rather than the objective.

This is not a new problem. It is recognizable in any long-horizon planning
process, in any multi-stage optimization, in any collaborative project where
the original intent must be transmitted across many hands over extended time.
What changes when search becomes abundant is not the structure of the problem
but its severity. When search is expensive, iterative processes are necessarily
short. The number of steps between the original specification and the final
output is limited by the cost of each step. There is less opportunity for the
objective to be progressively displaced by its proxy.

When search becomes cheap and iteration becomes rapid, the number of steps
that can occur between the original specification and any given intermediate
state grows dramatically. The original objective must now persist through far
more iterations. The accumulated displacement across many cheap steps can
exceed the displacement that would have occurred across a few expensive ones.

\begin{defn}
Let an objective $G$ be given at the outset of an inquiry. At each step $t$,
the agent's operative goal is a projection $G_t$ of $G$ onto the locally
visible context. A projection discards some aspects of $G$ while preserving
others. When the context is impoverished relative to $G$ --- when only a
small portion of the original objective is visible from the current intermediate
state --- the projection $G_t$ may differ substantially from $G$ while appearing
locally coherent. Over a sequence of steps, the cumulative displacement
$\|G - G_t\|$ may grow even when each individual transition $\|G_{t-1} - G_t\|$
is small.
\end{defn}

The failure mode this describes is distinctive. No individual step is obviously
wrong. Each local decision may appear reasonable given what is currently visible.
The trajectory as a whole, however, diverges from the original aim. The system
optimizes a shadow of the goal rather than the goal itself.

A simple illustration makes the structure concrete. Suppose the original
objective is to build a research platform that is flexible enough to accommodate
unanticipated future requirements. At each iteration, the locally visible
problem differs: in the first, the system is too slow; in the second, it is
too complex to explain to collaborators; in the third, it scores poorly on
a standard benchmark; in the fourth, it is difficult to deploy. Each optimization
is reasonable. Faster is better; simpler explanations are better; benchmark
scores matter; deployment ease matters. Yet the cumulative effect of optimizing
for speed, then simplicity, then benchmark score, then deployment ease is a
system that is specialized, well-documented for its current use, competitive
on current metrics, and easy to deploy in its current configuration --- but
no longer flexible. The original goal has disappeared not through any single
bad decision but through the accumulated displacement of many locally good ones.
The context visible at each step included the current problem and excluded the
original objective. Each projection was impoverished in the same direction.

Mechanisms that preserve goal continuity are therefore not merely useful
supplements to rapid generation. They are the structural counterpart without
which cheap search produces trajectories rather than progress. Such mechanisms
include: explicit review that compares intermediate states against original
objectives rather than against the immediately preceding state; deliberate
re-articulation of the objective at intervals throughout the inquiry; maintenance
of documentation that preserves the reasoning behind early commitments; and
periodic stepping back from local improvement to assess whether the overall
trajectory remains coherent.

None of these mechanisms are glamorous. None of them resemble the visible
productivity gains associated with rapid generation. But their value increases
precisely as search becomes cheaper. When generation is abundant and iteration
is rapid, the risk of cumulative displacement increases. The mechanisms that
counteract displacement become proportionally more important, not less.

% -----------------------------------------------------------
\section{Repair as Navigation}

The preceding sections have described several ways in which orientation, goal
continuity, and the structure of productive inquiry change when search becomes
abundant. A common thread connects them: activities that appear to be about
maintaining or recovering from failure are actually about preserving navigability
through time.

The dominant view of repair in complex systems treats it as cleanup. Something
has gone wrong; repair restores the system to correct functioning; inquiry
continues. On this view, a well-designed process would minimize the need for
repair. Repair is a cost, and its reduction is a goal.

An alternative view treats repair as the primary mechanism by which living
systems --- systems that persist and adapt through time rather than being
built once and completed --- maintain the capacity for future movement. On
this view, repair is not evidence of failure. It is the activity through which
a system stays navigable. A system that is never repaired is a system whose
trajectory has been fixed. A system that is regularly repaired is a system
that continues to have a trajectory.

\begin{principle}
Repair is not the correction of deviation from a fixed target. It is the
process by which an adaptive system maintains the conditions for continued
navigation. The value of repair is not that it restores a previous state but
that it preserves the possibility of future states.
\end{principle}

The distinction between these two views has practical consequences that become
clearer as search capacity increases. When generation is cheap and output is
abundant, a natural response is to regard repair as increasingly unnecessary.
If producing more output is nearly free, why invest in maintaining what has
already been produced? The cost-benefit calculation appears to favor generation
over repair.

This calculation neglects a dimension that does not appear in the immediate
cost structure: navigability. A system that has been extended without repair
accumulates structural commitments that were made in the context of earlier
states. Those commitments may no longer be appropriate given later developments.
They may be locally invisible --- no individual component of the system is
obviously wrong --- but they constrain the range of future modifications.
The cost of those constraints does not appear until future modifications are
attempted. At that point, the accumulated navigability debt becomes apparent.

The objection to rapidly generated, unrepaired systems is therefore not that
any given component is incorrect. It is that the future becomes difficult.
The system enters states from which many desirable future states are unreachable.
This is a reachability failure rather than a correctness failure. The system
is not wrong in the sense of failing to match a specification. It is wrong in
the sense of having closed off paths that might have led somewhere valuable.

Repair, in this light, is navigation. It is the activity through which the
agent maintains access to the future. Review, refactoring, constraint formation,
interface redesign, and explicit articulation of module boundaries are all
forms of repair. They appear inefficient from the perspective of immediate
generation. They are indispensable from the perspective of continued navigability.

The value of these activities is not constant. It increases as the rate of
generation increases. A system generated slowly accumulates navigability debt
slowly. A system generated rapidly accumulates navigability debt rapidly.
The faster the search, the more important the repair.

An instructive analogy appears in diffeomorphic image registration. The
objective of such a procedure is not merely to move one image into correspondence
with another. The transformation must preserve topology throughout the process:
no tearing, no folding, no destruction of the structural relationships that make
the final alignment meaningful~\cite{fireants2026}. A registration that reaches
the correct endpoint while introducing topological defects is regarded as
defective even if the final alignment appears locally successful. The reason
is that topology preservation constrains the admissible trajectory rather than
merely the endpoint. Quality depends not only on where the transformation
arrives but on what structural properties remain intact after arrival.

Repair in cognitive and technical systems plays an analogous role. Its purpose
is not to optimize the current state but to preserve the structural conditions
that keep future states reachable. A system that is locally correct but
topologically compromised --- one whose accumulated navigability debt has closed
off paths that cannot easily be reopened --- has failed in the same sense as
a registration that achieves correspondence while destroying the structure it
was meant to preserve.

% -----------------------------------------------------------
\section{The New Bottleneck}

The preceding analysis can be summarized as a sequence of displacements.

When search was expensive, the bottleneck in productive inquiry was the capacity
to generate candidates. More thinking, more design effort, more implementation
capacity --- these were the resources whose expansion would accelerate inquiry
most directly. Orientation was also required, but it was rarely the limiting
factor, because the cost of search naturally constrained the length of any
iterative process and thereby limited the opportunity for cumulative displacement.

When search becomes cheap, the bottleneck shifts. Generation is no longer
limiting. The capacity to orient --- to maintain coherent objectives through
extended iterations, to distinguish exploration from formalization, to preserve
goal continuity against cumulative displacement, to repair for navigability
rather than for immediate correctness --- becomes the resource whose scarcity
most constrains productive inquiry.

This shift has several implications.

First, activities that increase search capacity without attending to orientation
may produce less than their apparent output suggests. The visible quantity of
generated material increases. The proportion of that material that is genuinely
useful --- that has been selected, evaluated, and integrated in relation to a
well-maintained objective --- may not increase proportionally and may in some
cases decrease.

Second, the contribution of human judgment to any collaborative inquiry between
a person and a generation system is not symmetric with the system's contribution.
The system contributes search. The human contributes orientation. These are not
interchangeable. A more powerful search system does not reduce the need for
orientation; it increases it, because a larger explored space requires more
careful navigation. The appropriate response to more capable generation is not
less human involvement in the inquiry but differently structured human involvement
--- concentrated in the orientation activities that search cannot perform.

Third, the skills most valuable in productive inquiry have shifted. When search
was expensive, implementation skill --- the ability to generate high-quality
outputs efficiently --- was central. When search becomes abundant, implementation
skill remains necessary but is no longer distinctive. What becomes distinctive
is the capacity for orientation: the ability to maintain coherent objectives
through extended iteration, to recognize when exploration should give way to
commitment, to evaluate trajectories rather than states, and to repair for
navigability rather than for immediate correctness. These are not new skills.
They are skills that good inquiry has always required. What has changed is their
relative value.

Fourth, the appropriate structure of productive work changes. When generation
was expensive, it made sense to minimize the number of iterations by investing
heavily in specification before implementation. When generation is cheap, this
logic reverses: it becomes sensible to invest less in prior specification and
more in exploration, using abundant generation to examine the territory before
committing to an architecture. But exploiting this reversal requires maintaining
the discipline of orientation throughout the exploration phase, so that the
insights gained from exploration are not lost when formalization eventually occurs.

% -----------------------------------------------------------
\section{Formal Interpretation}

The preceding account has been developed in ordinary language, with the aim of
making the central asymmetry accessible without presupposing any particular
formal apparatus. It may nevertheless be useful to indicate briefly how the
observations can be expressed in a more precise vocabulary.

The concept that provides the most direct formal grip on the phenomena described
here is \emph{projection collapse}.

Let $G$ denote the goal defined at the outset of an inquiry. At each stage $t$
of an iterative process, the agent does not have direct access to $G$. The
agent has access to the current intermediate state $S_t$ and to some context
$C_t$ capturing what is locally visible from that state. The operative goal at
stage $t$ is a projection $\pi(G, C_t)$: the image of $G$ restricted to the
locally visible context. Projection is necessarily lossy. When $C_t$ is rich,
$\pi(G, C_t)$ is faithful and optimization serves $G$ well. When $C_t$ is
impoverished, the projection is distorted and optimization relative to
$\pi(G, C_t)$ may diverge substantially from $G$.

\begin{defn}[Projection Collapse]
A sequence of projections $\{\pi(G, C_t)\}_{t=0}^{T}$ undergoes projection
collapse if the cumulative displacement $d(G, \pi(G, C_T))$ grows substantially
over the sequence, even when each individual step displacement
$d(\pi(G, C_{t-1}), \pi(G, C_t))$ is small. Projection collapse occurs when
iterative local optimization produces global divergence from the original objective.
\end{defn}

Projection collapse is the formal analog of cumulative goal displacement. The
trajectory appears locally coherent: each step is reasonable given what is
currently visible. The sequence as a whole, however, moves away from the
original aim, because each step optimizes a projection that has already lost
some information about $G$.

The phenomena described in earlier sections map directly onto this structure.
Exploration before formalization is the practice of ensuring $C_t$ remains
rich before commitment narrows the accessible context. Goal continuity
mechanisms --- review, re-articulation, documentation of early reasoning ---
are practices for maintaining a representation of $G$ that is not entirely
determined by $C_t$, periodically re-anchoring the operative goal against
accumulated drift. Repair as navigation is the practice of preventing structural
commitments from impoverishing the context to the point where future states
become unreachable.

Orientation, in these terms, is the capacity to maintain a faithful correspondence
between $\pi(G, C_t)$ and $G$ across an extended sequence --- to detect when
the projection has become impoverished, and to act on that detection. This
capacity cannot be delegated to any process that operates only within the local
context. It is the irreducible contribution of the inquirer to any inquiry in
which generation has been substantially externalized. As search scales, that
contribution does not diminish. It becomes the limiting resource.

% -----------------------------------------------------------
\section*{Conclusion}

The collapse of generation costs has expanded search capacity in ways that are
genuine and substantial. More of possibility space is accessible to individual
inquiry than at any previous point. More alternatives can be examined before
commitment. More candidate structures can be held simultaneously in play. These
are real enlargements of what inquiry can accomplish.

What has not expanded is orientation capacity. The higher levels of
orientation --- goal maintenance, trajectory assessment, and goal revision
--- remain tied to the cognitive and institutional resources of the inquiring
agent. They do not scale computationally. They become relatively more scarce
as search becomes relatively more abundant.

The consequence is a shift in what matters most. When generation was the
bottleneck, generating well was the primary challenge. When generation becomes
abundant, remaining oriented becomes the primary challenge. Activities that
appear to be about maintenance, review, repair, and explicit goal articulation
are not inefficiencies to be minimized. They are the mechanisms through which
productive inquiry remains productive as the scale of search increases.

This principle does not depend on any particular technology or moment. It applies
wherever the capacity to generate candidates substantially exceeds the capacity
to evaluate trajectories. It is as relevant to scientific programs that have
grown faster than their conceptual frameworks, to institutions that have expanded
faster than their governance structures, and to bodies of knowledge that have
accumulated faster than their organizing principles, as it is to any other
domain where search and orientation have become uncoupled.

The central problem of an age of abundant search is not how to generate
possibilities. It is how to navigate among them without losing the goal that
made them worth exploring.

% -----------------------------------------------------------
\bigskip
\begin{center}
  \rule{0.4\textwidth}{0.4pt}
\end{center}

